%% ============================================================
%% sec01_intro.tex — §1 引言 (Introduction)
%% 三线建设 Exit-DDD 顶刊研究 (中国工业经济 / 经济研究 范式)
%% ============================================================
\section{引言}\label{sec:intro}

中国经济体制的"产业政策实验"在 20 世纪后半叶为研究者提供了若干被学术界反复利用的准自然实验来源：1980 年代的家庭联产承包责任制、1990 年代的国有企业"抓大放小"改制、2001 年加入 WTO 后的贸易自由化冲击，以及 2013 年以来的"一带一路"倡议。然而，1964--1980 年间贯穿川黔陕甘等十三省、累计投资逾 2{,}000 亿元、布局约 1{,}100 家大中型军工与配套企业的\textbf{三线建设}，因其涉密程度高、空间分布极不均衡、且牵涉"进、出、迁、并、停"五类政策动作的叠加，长期游离于主流经济学实证文献的视野之外。\citet{FanZou2021} 关于"产业政策进入端"的开创性研究表明，国家主导的三线建设重工业空间集聚对当地经济具有显著的长期因果效应；但\textbf{产业政策的"退出端"——即国有三线企业从山沟"脱险搬迁"到城市、关停并转所释放出的本地沉淀人力资本与供应链网络——尚未被定量识别}。这是本文要切入的第一个研究空白。

现有文献在两条线索上互为补充但缺乏整合。第一条线索聚焦\textbf{国有企业的"关停并转"行为}。\citet{hsieh2009}、\citet{brandt2013}、\citet{li2017} 在中国制造业生产率分解的框架下，识别了 1990 年代国有企业"抓大放小"对全要素生产率的提升效应；\citet{白重恩2016} 与 \citet{聂辉华2016} 进一步从政治锦标赛视角解释了"为何落后国企不被允许破产"。但这些研究多以全国或省级面板为单位，\textbf{难以识别"同一城市-行业内部，国有三线企业撤并"对非国有制造业新生企业的因果效应}。第二条线索聚焦\textbf{产业集聚与扩散}。\citet{marshall1920}、\citet{krugman1991}、\citet{combes2000} 等指出，知识溢出与劳动力蓄水池构成集聚经济两大微观基础；\citet{Combes2010} 与 \citet{marcon2003} 用法国植物专利数据证实知识溢出存在明显的地理衰减与行业相关性。\citet{FanZou2021}、\citet{陆铭2016} 也在国内市场分割与高铁通达性的语境下，重申了"地方化知识溢出"对民营经济的重要性。国内三线建设经济学实证方面，\citet{徐现祥2017} 以"中央计划体制的退出"为核心视角开创了国内三线撤并端的研究先河；\citet{刘冲2020} 则利用"三线建设与西部地区经济发展"的微观地理识别，将三线进入端的因果识别精度推进到"县-产业"维度。然而，\textbf{这些文献普遍将"国有撤并"视作给定的外生冲击，没有追问——如果大树拔走后，本地民营经济是同步凋零、还是反而被留下的"看林人"激活？} 我们利用 1986 年开始的三线企业"调整改造"这一外生政策窗口，把"三线企业撤并"作为产业政策从"进入"切换到"退出"的反向准自然实验，提供对这一空白问题的因果回答。

本文选取\textbf{四川、贵州两省 92 家最具代表性的三线大中型企业}为种子样本，构造 1980--2017 年县级 × 二位数制造业 (CIC-2) × 年的三维非平衡面板，采用\textbf{三重差分 (DDD) + 微观地理长差分 / 面板工具变量 (IV) + Chen \& Roth (2024) 零值稳健检验}的干净识别策略。研究发现，三线企业撤并对原驻留县域相同与上下游关联二位数行业内的非国有制造业企业诞生率具有\textbf{显著正向且经济上可观的促进效应}：以非国有制造业企业新登记数衡量，处理组在撤并后的 5 年内新登记企业密度较对照组上升约 8\%--12\%（IV-DDD 估计值），而以"销售收入 / 资产规模"等增量边际指标衡量，效应在撤并后第 3 年达峰并持续约 7 年。异质性分析进一步表明，撤并的"创业孵化"效应在\textbf{具备通用目的技术 (General Purpose Technologies, GPTs) 的电子元器件、通用机械、精密仪表行业}显著高于\textbf{资产专用性较强的重装火工、航空发动机、专用兵器}等行业，验证了知识可迁移性作为溢出渠道的核心条件。\citet{bresnahan1995}、\citet{helpman1998} 强调的 GPTs 范式由此获得来自计划经济体制退出实验的微观证据。

本文的边际贡献可概括为以下三个层面。\textbf{第一，识别策略层面}：将主流文献长期聚焦的"产业政策进入端"（如 \citet{FanZou2021} 关于"三线建设"对当地经济长期发展的因果识别）拓展到"退出端"，构造了\textbf{国有三线企业"撤并"对地方民营经济溢出效应的第一份三维 DDD 识别}。相较 \citet{徐现祥2017} 用省际面板考察"中央计划体制退出"、\citet{刘冲2020} 用县-产业面板考察"三线进入端"的策略，本文利用"县-行业-年"三维面板吸收了城市-行业双维固定效应，并通过微观地理 IV (HasTF × Post) 缓解了"地方官员选择性推动撤并"的内生性。\textbf{第二，机制层面}：本文发现"撤并"并非简单的"大树拔走、生态凋零"，而是触发了\textbf{沉没人力资本 (Hysteresis in Sunk Human Capital) 与供应链网络粘性}两类反向力量——前者体现在"厂办技校/子弟校"毕业生的本地创业溢价上，后者表现为原配套零部件厂商被本地民营接盘后的"接链效应" (link-takeover effect)。这一发现与 \citet{hausmann2003} 关于产业复杂度、\citet{hidalgo2009} 关于经济复杂性的预言互为印证：在具备 GPTs 基础的多样化工业生态中，撤并反而通过"被留下的看林人"激活了新一轮的本土企业家精神 (entrepreneurial spawning)。\textbf{第三，方法论层面}：本文为"国家主导的产业空间撤并"研究提供了一套可复制的识别模板——它包括 (i) 一个由 \textbf{92 家真实三线企业}组成的种子库 (B0\_third\_front\_seed\_db.csv)，(ii) 三维面板 DDD 配合微观地形 (HasTF) IV 的双重识别，(iii) Borusyak et al. (2024) \texttt{did\_imputation} 与 Callaway \& Sant'Anna (2021) \texttt{csdid} 提供的异质性稳健 DID 双重背书，(iv) \citet{chen2024} 零值效应 (zero-event) 检验对"无效假说"的可证伪性 (falsifiability) 加固。这一组合识别模板可被推广到 1990 年代的"国企抓大放小"、2008 年汶川地震后的"产业重建"以及 2010 年代东北"老工业基地振兴"等类似政策窗口。

本文其余部分安排如下：第~\ref{sec:bg} 节梳理三线建设的制度演进与本文研究背景，并提出三个可检验的机制假说；第~\ref{sec:strategy} 节交代识别策略、数据来源与关键变量构造；第~\ref{sec:main} 节报告主回归与异质性分析结果；第~\ref{sec:robust} 节提供一系列稳健性检验；第~\ref{sec:conclusion} 节对三大机制假说分别做证据穿透；第~\ref{sec:conclusion} 节总结结论并讨论政策启示。
